\documentclass[letterpaper]{article}

% AAAI 2026 style.
% Use [submission] if you want the anonymous/submission version.
% For a course report with visible names, try without [submission].
\usepackage{aaai2026}

% Required / standard AAAI packages
\usepackage{times}
\usepackage{helvet}
\usepackage{courier}
\usepackage[hyphens]{url}
\usepackage{graphicx}
\usepackage{natbib}
\usepackage{caption}

% Useful packages
\usepackage{amsmath}
\usepackage{amssymb}
\usepackage{booktabs}
\usepackage{multirow}

% Optional: algorithms
\usepackage{algorithm}
\usepackage{algorithmic}

\urlstyle{rm}
\def\UrlFont{\rm}
\frenchspacing

% US Letter page size, as required by AAAI
\setlength{\pdfpagewidth}{8.5in}
\setlength{\pdfpageheight}{11in}

% No section numbering, AAAI-style
\setcounter{secnumdepth}{0}

\pdfinfo{
/TemplateVersion (2026.1)
}

\title{Rule-Based and LLM-Assisted Evaluation of Counterfactual Explanations for Income Prediction}

\author{
    AFRA Lama,
    MORTADA Daniel,
    RIABICHEFF Ivan,
    DESOIL Théo
}

\affiliations{
    Universit\'e Libre de Bruxelles\\
    INFO-H512 Project\\
    \{lama.afra, daniel.mortada, ivan.riabicheff, theo.desoil\}@ulb.be
}

\begin{document}

\maketitle
\thispagestyle{plain}
\pagestyle{plain}

\begin{abstract}
This project studies the evaluation of counterfactual explanations for income
prediction. We use DiCE to generate counterfactuals on the Adult Income dataset
and evaluate them with a deterministic metrics-only reference system. We then
test whether a calibrated single LLM can approximate this reference and whether
a second LLM call can provide useful human-readable explanations. The calibrated
single LLM reaches high issue-label agreement with the reference, while the
explainability layer adds qualitative value without changing the fixed verdicts.
\end{abstract}

\section{Introduction and Problem Formulation}
\noindent
This project connects two seminar themes: \textbf{Counterfactuals for
Explainability and Fairness in Machine Learning} and \textbf{Orchestrate,
Secure \& Scale Enterprise GenAI Architecture}. It is related to the first
seminar because we study counterfactual explanations for a binary income
prediction model. Given a person
whose income is predicted to be below \$50K, we use DiCE to generate
counterfactual examples that would change the model prediction to above \$50K.
The central question is not only whether these counterfactuals flip the model
output, but whether they are good explanations. In the context of explainable
AI, a useful counterfactual should be close to the original input, plausible,
actionable, and sparse, meaning that it should change only a small number of
features.\\
\newline
The project is also related to \textbf{Orchestrate, Secure \& Scale Enterprise
GenAI Architecture} because we implement an LLM-assisted evaluation pipeline.
The seminar discussed generative AI workflows, orchestration of AI services,
hallucination control, auditability, human validation, and trustworthy LLM
systems. Our final pipeline follows this perspective by using a calibrated
single-LLM evaluator and a separate explanation layer around a deterministic
reference system.\\
\newline
The main research question addressed in this report is:

\begin{quote}
    \emph{Can a calibrated single LLM approximate a deterministic metrics-only
    reference for counterfactual quality evaluation, and can an explanation
    layer add useful human-facing interpretation without changing the verdict?}
\end{quote}
\noindent
This question is relevant because counterfactual generation methods such as DiCE
typically select counterfactuals using deterministic criteria. These include
\textbf{validity}, which checks whether the generated counterfactual flips the
model prediction; \textbf{proximity}, which favors small changes from the
original input; \textbf{sparsity}, which favors changing few features; and
\textbf{feasibility}, which can restrict which features are allowed to change.
These metrics are necessary, but they do not fully capture whether a proposed
change is realistic, fair, causally coherent, or convincing from a human
perspective. We therefore test whether LLM-based evaluators provide additional
useful assessment of the counterfactuals generated by DiCE. Since no external
ground-truth labels exist for the quality of these counterfactual explanations,
we use a deterministic metrics-only evaluator as a rule-based reference
standard. This reference encodes our project-specific human judgment through
explicit heuristics and thresholds over validity, sparsity, proximity,
confidence, and plausibility checks. We then evaluate whether a single LLM can
substitute for this reference at the issue-label level and whether an optional
explanation layer makes the verdicts easier to inspect.\\
\newline
We evaluate this question on the OpenML Adult Income dataset, where the task is
to predict whether a person's income is above or below \$50K.

% The goal of this project is to study in depth a state-of-the-art AI method
% presented in the seminar. We focus on the topic of [seminar topic], and more
% specifically on [selected method].

% The main research question addressed in this report is:

% \begin{quote}
%     \emph{[Write your precise research question here.]}
% \end{quote}

% This question is relevant because [motivation]. In this work, we [reproduce /
% apply / extend / compare] the method by implementing it and evaluating it on
% [dataset or problem setting].

% Our contributions are:
% \begin{itemize}
%     \item a targeted literature review of [research area];
%     \item an implementation of [selected method];
%     \item an experimental evaluation on [dataset/task];
%     \item a critical analysis of the method's strengths, assumptions, and limitations.
% \end{itemize}

\section{Related Work}

\subsection{Taxonomy of the Research Area}

Machine learning is increasingly used in high-impact domains such as healthcare,
finance, and banking. This makes explainability important, because many useful
models are difficult to interpret directly. The research area can be divided
into three relevant families of methods:
\begin{itemize}
    \item \textbf{Intrinsically interpretable models:} methods such as linear
    models or decision trees are easier to inspect because their decision
    process is relatively transparent. However, they may be less flexible than
    more complex models.
    \item \textbf{Post-hoc counterfactual explanations:} these methods explain a
    model decision by asking what minimal changes to the input would have led to
    a different prediction. They are useful because they can explain individual
    decisions without requiring access to the internal structure of the model.
    \item \textbf{Counterfactual quality evaluation:} generated
    counterfactuals are usually assessed with metrics such as validity,
    proximity, sparsity, diversity, and feasibility. These metrics are useful
    but incomplete, because they do not always capture whether a recommendation
    is realistic, fair, or actionable in the real world.
\end{itemize}
\noindent
This taxonomy positions our project in the third family. We do not only generate
counterfactual explanations; we study how their quality should be evaluated. In
particular, we use deterministic quality metrics and heuristics as a stable
reference, then test whether an LLM can reproduce and explain that reference.

\subsection{Seminal and Recent Works}
\noindent
Counterfactual explanations were introduced as a way to explain automated
decisions without opening the black box of the underlying model
\citep{wachter2017counterfactual}. Instead of describing all internal model
mechanisms, a counterfactual explanation answers a local question: what would
need to change in this input for the model to produce a different outcome? For a
classification task, this means showing how an unfavorable prediction could be
changed into a favorable one.\\
\newline
Later work extended this idea by emphasizing diversity and actionability, rather
than validity alone. DiCE generates multiple diverse counterfactual alternatives
for a given instance, giving users several possible paths toward the desired
outcome instead of a single potentially unrealistic modification
\citep{mothilal2020explaining}. This is the counterfactual generation framework
used in our project.\\
\newline
However, DiCE-style evaluation still relies heavily on numerical criteria.
Validity checks whether the model prediction flips; proximity measures how close
the counterfactual is to the original instance; sparsity measures how many
features change; diversity measures variation among generated alternatives; and
feasibility constrains which features may be modified. These criteria are
necessary, but they can miss important semantic failures. For example, a
counterfactual may be close in feature space while still proposing an impossible
change, such as decreasing a person's age from 30 to 25. Similarly, numerical
metrics may not fully capture fairness, realism, or causal coherence.\\
\newline
Our project addresses this gap with a deterministic reference evaluator and an
LLM-assisted explanation layer. Because there is no external ground truth for
these quality labels, our metrics-only evaluator acts as the reference standard:
it encodes human judgment through explicit rules over metrics and
issue-specific plausibility checks. The LLM is then assessed by how well it
reproduces this reference and whether it can explain the result in clearer
language. To make the audit operational, we define a small issue taxonomy
covering implausible time-dependent changes, extreme working hours,
unactionable capital shifts, too many simultaneous changes, and fragile
counterfactuals; the next section describes how these labels are used in the
evaluation pipeline.

\section{Method}

The original predictive task is binary classification on the Adult Income
dataset, where the model predicts whether an individual's income is at most or
above \$50K. We train a Logistic Regression model inside an \texttt{sklearn} pipeline. 
This model is used for counterfactual generation because DiCE requires reliable class-probability
outputs that can guide its optimization procedure.\\
\newline
Before generating counterfactuals with DiCE, we define a feature policy that
freezes attributes that should not be modified: \texttt{fnlwgt}, marital
status, relationship, race, sex, and native country. Counterfactuals are
therefore not allowed to assign different values to these features, since doing
so would imply unfair or unrealistic demands on the individual.\\
\newline
After training the predictive model and applying this feature policy, we
generate counterfactual explanations with DiCE. For each individual with an
unfavorable prediction, DiCE produces multiple candidate counterfactuals that
aim to move the prediction to the favorable class. We keep the DiCE generation
parameters at their default values and use the feature policy only to restrict
which attributes may change.\\
\newline
For each generated counterfactual, we compute a first set of standard numerical
quality metrics:
\begin{itemize}
    \item \textbf{Validity:} whether the counterfactual reaches the favorable
    class.
    \item \textbf{Proximity:} how close the counterfactual remains to the
    original instance.
    \item \textbf{Sparsity:} how many features are changed.
    \item \textbf{Diversity:} how different the generated counterfactuals are
    from each other.
    \item \textbf{Counterfactual confidence:} the model probability assigned to
    the favorable class for the counterfactual.
\end{itemize}
\noindent
These metrics are useful for ranking and filtering counterfactuals, and they
are close to the criteria commonly used in DiCE-style generation. However, they
remain purely numerical. They do not directly represent whether a
counterfactual is plausible, actionable, or convincing from a human perspective.
We therefore add a five-label issue taxonomy that provides a second evaluation
layer. The taxonomy defines a fixed vocabulary of possible counterfactual
failures:
\begin{itemize}
    \item \texttt{implausible\_time\_dependent\_change}: age decreases,
    education decreases, education increases without sufficient age increase,
    or age/education becomes non-integer.
    \item \texttt{extreme\_working\_hours}: weekly working hours become
    dangerously high, implausibly low, or change by an unrealistic amount.
    \item \texttt{unactionable\_capital\_shift}: capital gain or capital loss
    changes in a financially unrealistic way.
    \item \texttt{too\_many\_changes}: too many actionable features are changed
    at the same time.
    \item \texttt{fragile\_counterfactual}: the counterfactual only barely
    reaches the favorable class, with confidence close to the decision
    threshold.
\end{itemize}
\noindent
Each label corresponds to a specific failure mode in the generated
counterfactual. The previous label
\texttt{inconsistent\_work\_profile} was removed from the active taxonomy
because the deterministic rules did not detect it reliably and LLM outputs
tended to overuse it. This reduces noisy labels and keeps all evaluation
strategies aligned with the same vocabulary.\\
\newline
We operationalize the taxonomy through deterministic heuristic rules. These
rules encode human-designed plausibility checks over the original instance, the
counterfactual, and the numerical metrics. For example, they flag decreases in
age, unrealistic changes in education or working hours, large unexplained
capital shifts, excessive simultaneous modifications, and low-margin favorable
predictions. The heuristic module returns the set of issue labels detected for
each counterfactual. We also compute metric warnings, including
\texttt{count\_diversity\_low}, which flags counterfactual sets with low
diversity. These warnings inform severity but are not scored taxonomy labels.\\
\newline
Finally, we define a metrics-only deterministic evaluator. This evaluator does
not use an LLM. It receives as input the generated counterfactuals, their
standard numerical metrics, and the heuristic issue flags derived from the issue
taxonomy. It outputs two judgments for each counterfactual:
\begin{itemize}
    \item \textbf{Issue severity:} \texttt{low}, \texttt{medium}, or
    \texttt{high}, depending on the number and type of detected issues.
    \item \textbf{Overall assessment:} \texttt{fair}, \texttt{ambiguous}, or
    \texttt{unfair}, summarizing whether the counterfactual is acceptable as an
    explanation.
\end{itemize}
\noindent
This deterministic evaluator acts as the rule-based reference against which the
LLM-based evaluation strategies are compared.\\
\newline
\textbf{Single-LLM substitution evaluator.} The main LLM-based system is a
single evaluator prompted to approximate the metrics-only reference. It receives
the original profile, generated counterfactuals, feature changes, numerical
metrics, heuristic issue evidence, and taxonomy instructions. The prompt is
calibrated to:
\begin{itemize}
    \item approximate the metrics-only reference rather than invent a new
    evaluation policy;
    \item start from the heuristic issue union
    \texttt{flagged\_issues\_union};
    \item keep labels supported by deterministic heuristic evidence;
    \item avoid unsupported labels outside the active taxonomy;
    \item preserve \texttt{fragile\_counterfactual} when it is present in
    case-level or counterfactual-level heuristic flags.
\end{itemize}
\noindent
The output is a structured verdict with an overall assessment
(\texttt{fair}, \texttt{ambiguous}, or \texttt{unfair}), detected issue labels,
severity, confidence, and recommended action.\\
\newline
\textbf{Explanation layer.} In the final version, explainability is added as a
second LLM call after the verdict is fixed. This layer receives the fixed
single-LLM verdict and compact case evidence, then produces an
\texttt{expert\_explanation} for a non-expert reader. It is explicitly
forbidden from changing issue labels, severity, assessment, confidence, or the
recommended action. This separates the substitution task from the human-facing
interpretation task.\\
\newline
\textbf{Earlier multi-agent system.} A multi-agent AutoGen debate with
Prosecutor, Defense attorney, Expert witness, and Judge was implemented during
the project. In the final Phase 2 framing, however, the central evaluation is
the single-LLM substitution and explanation pipeline, because the goal is to
measure approximation of the deterministic reference rather than to prove that
debate is more correct.

% We now describe [selected method]. The method aims to [goal of method].
% At a high level, it works by [main algorithmic idea].

% Given an input $x$, the method computes [output] by solving/estimating:

% \begin{equation}
%     \hat{y} = f_\theta(x),
% \end{equation}

% where $f_\theta$ denotes [model/method]. The key steps are:

% \begin{enumerate}
%     \item [Step 1];
%     \item [Step 2];
%     \item [Step 3].
% \end{enumerate}

% In our implementation, we follow [paper/source] with the following adaptations:
% [describe simplifications, modifications, or implementation choices].

\section{Experimental Setup}

\subsection{Dataset}

We evaluate the pipeline on the Adult Income dataset. The task is binary
income prediction: each individual is classified as earning at most \$50K or
more than \$50K per year.

\begin{table}[H]
\centering
\small
\begin{tabular}{p{0.36\columnwidth}p{0.54\columnwidth}}
\toprule
Property & Description \\
\midrule
Dataset & Adult Income \\
Instances & 48,842 \\
Input features & 14 demographic, educational, occupational, financial, and
working-time attributes \\
Prediction task & Binary income classification \\
Unfavorable class & Predicted income $\leq$ \$50K \\
Favorable class & Predicted income $>$ \$50K \\
\bottomrule
\end{tabular}
\caption{Dataset used in the experiments.}
\label{tab:dataset}
\end{table}

\noindent
The main features include age, education, occupation, capital gain, capital
loss, hours worked per week, and native country. For the evaluation stage:
\begin{itemize}
    \item we select 10 individuals predicted in the unfavorable class;
    \item DiCE generates multiple counterfactuals for each selected individual;
    \item DiCE generation uses its default parameters, with the feature policy
    described in the Method section.
\end{itemize}

\subsection{Baselines}

We compare the LLM systems against the metrics-only evaluator. This evaluator is
not an external ground truth, but it is the deterministic reference system used
for substitution scoring.
\begin{itemize}
    \item \textbf{Metrics-only reference:} deterministic evaluator based on
    numerical metrics and heuristic issue flags.
    \item \textbf{Uncalibrated single LLM:} one LLM produces a structured
    verdict from the same case evidence.
    \item \textbf{Calibrated single LLM:} the single-LLM prompt is constrained
    to approximate the metrics-only reference and preserve heuristic labels.
    \item \textbf{Calibrated single LLM with explanation:} the calibrated
    verdict is followed by a second LLM call that explains the fixed verdict.
\end{itemize}

\subsection{Metrics}

Evaluation focuses on agreement with the metrics-only reference:
\begin{itemize}
    \item \textbf{Detection rate:} percentage of taxonomy issue labels from the
    metrics-only reference that are also identified by the evaluated system.
    \item \textbf{False-positive rate:} percentage of cases where the evaluated
    system adds issue labels absent from the reference.
    \item \textbf{Exact match rate:} percentage of cases for which the complete
    predicted issue-label set exactly matches the metrics-only reference set.
\end{itemize}

\subsection{Implementation Details}

The implementation is written in Python. Counterfactuals are generated with
DiCE, and the LLM-based evaluators are implemented with AutoGen. The final
Phase 2 experiments use Groq with \texttt{llama-3.1-8b-instant}.

\begin{table}[H]
\centering
\small
\begin{tabular}{lcc}
\toprule
Setting & Calibrated LLM & LLM + explanation \\
\midrule
Framework & AutoGen & AutoGen \\
Provider & Groq & Groq \\
Model & \multicolumn{2}{c}{\texttt{llama-3.1-8b-instant}} \\
Temperature & \multicolumn{2}{c}{0.2} \\
Max tokens & 400 & 700 \\
\bottomrule
\end{tabular}
\caption{LLM configuration used in the final Phase 2 experiments.}
\label{tab:llm-config}
\end{table}

\noindent
The temperature is set to 0.2 to keep outputs conservative while avoiding fully
deterministic generation. The final explainability run uses a larger completion
budget because it adds a second-stage human-readable explanation. Earlier
multi-agent experiments used a smaller 250-token budget to stay under Groq's
token limits, but the final Phase 2 evaluation focuses on the single-LLM
substitution and explanation pipeline.

\section{Results and Analysis}

Table~\ref{tab:results} summarizes the main experimental results.

\begin{table}[H]
\centering
\small
\begin{tabular}{lccc}
\toprule
System & Detect. & FP & Exact \\
\midrule
Uncalibrated single LLM & 84.0\% & 0.0\% & 60.0\% \\
Calibrated single LLM & 96.0\% & 0.0\% & 90.0\% \\
LLM + explanation layer & 96.0\% & 10.0\% & 80.0\% \\
\bottomrule
\end{tabular}
\caption{Substitution scores against the metrics-only reference.}
\label{tab:results}
\end{table}

Prompt calibration substantially improves issue-label substitution. Detection
increases from 84.0\% to 96.0\%, and exact match increases from 60.0\% to
90.0\%. This shows that the LLM performs best when it is instructed to follow
the deterministic heuristic evidence rather than act as a free-form evaluator.\\
\newline
The final explainability run keeps the same 96.0\% detection rate but has lower
exact match, 80.0\%, and a 10.0\% false-positive rate. This is interpreted as
normal LLM variability in a fresh run, not as the explanation layer changing the
labels: the explanation is generated only after the verdict is fixed.

\section{Discussion}

The results support the Phase 2 framing of the project. The deterministic
metrics-only evaluator is the stable backbone: it defines the reference issue
labels and keeps the evaluation auditable. The calibrated single LLM can
approximate this reference well at the issue-label level, but only when the
prompt explicitly tells it to preserve deterministic heuristic evidence. This
suggests that the LLM is more reliable as a structured substituter than as an
independent judge.\\
\newline
The explanation layer provides the clearest added value. It does not improve the
quantitative substitution scores, and it is not allowed to change the verdict.
Instead, it translates the fixed structured assessment into language that is
easier for a reader to understand. This is useful because the metrics-only
reference is precise and reproducible, but its raw issue labels and thresholds
are not always self-explanatory.\\
\newline
The main limitation is that the reference itself is rule-based. It reflects our
designed heuristics, not an external human-labeled ground truth. In addition,
LLM outputs remain variable across runs, as shown by the difference between the
best calibrated run and the final explainability run. Future work could test a
larger model, add human annotations, or calibrate the metric-warning thresholds
such as \texttt{count\_diversity\_low}.

\section{Conclusion}

This report studied counterfactual explanation evaluation for income prediction.
We generated counterfactuals with DiCE, evaluated them with a deterministic
metrics-only reference, and tested whether a single LLM could approximate and
explain that reference. The final conclusion is that deterministic rules remain
the evaluation backbone, while the calibrated LLM is useful as a low-cost
substitution and explanation layer. Its strongest contribution is not superior
correctness, but clearer human-facing interpretation of an auditable rule-based
verdict.

\bibliography{references}

\end{document}
